% !TeX root = RJwrapper.tex
\section*{Supplementary Materials}
\section{Additional real data results}

\begin{figure}[!ht]
\centering
\begin{subfigure}[b]{0.48\textwidth}
    \centering
    \includegraphics[width=\textwidth]{real data application figures/new_real_data_figures/randomforest credit.pdf}
    \caption{Random forest}
    \label{fig:german-rf}
\end{subfigure}
\hfill
\begin{subfigure}[b]{0.48\textwidth}
    \centering
    \includegraphics[width=\textwidth]{real data application figures/new_real_data_figures/SVM credit.pdf}
    \caption{SVM}
    \label{fig:german-svm}
\end{subfigure}

\caption{\footnotesize Distribution of the approximate errors for South German credit data analysis under different base classification methods.}
\label{fig:german_compare}
\end{figure}


\begin{table}[h!]
\centering
\resizebox{\textwidth}{!}{%
\begin{tabular}{c !{\vrule width 1.2pt} c !{\vrule width 1.2pt}  c c c c !{\vrule width 1.2pt} c c c c !{\vrule width 1.2pt} c}
\Xhline{1.2pt}
Base method & Paradigm 

& $R_{1\star}$ & $R_{2\star}$ & $R_{3\star}$ & $R_{4\star}$
& $V_1$ & $V_2$ & $V_3$ & $V_4$ & $R_{\text{overall}}$\\
\Xhline{1.2pt}

\multirow{2}{*}{Random forest}
& Classical 
 & 0.095 & 0.158 & 0.084 & 0.112 
& 0.020 & 0.280 & 0.010 & 0.080 & 0.572 \\
& \textcolor{black}{H--NP}  
 & 0.132 & 0.077 & 0.045 & 0.063 
& 0.110 & 0.050 & 0.000 & 0.020 & 0.617\\

\Xhline{1.2pt}

\multirow{2}{*}{SVM}
& Classical 
 & 0.449 & 0.234 & 0.113 & 0.119 
& 1.000 & 0.680 & 0.070 & 0.030 & 0.618 \\
& \textcolor{black}{H--NP}  
 & 0.140 & 0.061 & 0.025 & 0.028 
& 0.150 & 0.020 & 0.000 & 0.000 & 0.654\\

\Xhline{1.2pt}

\end{tabular}
}
\caption{Average overall classification error $R_{\text{overall}}$, average under-classification errors $R_{i\star}$'s, along with their corresponding violation rates $V_i$'s for the German credit data analysis.}
\label{tab:german_credit_base_methods}
\end{table}


\section{Additional simulation results}\label{sec:extra-sim}

% \section{Extra examples}








% \subsection{Air Quality Classification}
% We show \texttt{hnp\_umbrella()} based on multinomial logistic regression, multi-class support vector machine and random forest classifiers is able to control under-classification errors. Firstly, we trained 3 base classifiers using the following method:

% We implemented multinomial logistic regression using the \texttt{nnet::multinom(y \~{} ., data = dataset, trace = FALSE)} function from the \texttt{nnet} package. Class probabilities required for the H--NP construction were estimated using \texttt{predict(fit, Xnew, type = "probs")}.

% The SVM base learner was implemented via the \texttt{e1071::svm(y \~{} ., data = dataset, scale = TRUE, probability = TRUE)} function. Unless otherwise specified, the model utilized the default radial basis function (RBF) kernel, with the cost and $\gamma$ parameters kept at their default values without additional tuning. Test-set predictions were obtained via \texttt{predict(fit, Xnew)}, and the required probabilities were extracted using \texttt{attr(predict(fit, Xnew, probability = TRUE), "probabilities")}.

% We utilized the \texttt{randomForest::randomForest(x = x, y = y, ntree = 100, mtry = 2)} function in classification mode. For performance evaluation, test-set predictions were obtained via \texttt{predict(fit, Xnew, type = "class")}, while class probabilities for the H--NP construction were estimated using \texttt{predict(fit, Xnew, type = "prob")}.

% We demonstrate how to use \texttt{hnp\_umbrella()} to control under-classification errors in an air-quality classification task at prespecified control levels. The dataset is available at \url{https://www.kaggle.com/datasets/mujtabamatin/air-quality-and-pollution-assessment/data}. The dataset contains measurements of air-quality indicators such as PM 2.5 and temperature, together with a corresponding air-quality label. The label contains 4 levels from bad to good: Hazardous, Poor, Moderate, Good. With this unbalanced dataset, we do the corresponding mapping:
% \begin{table}[htbp]
% \centering
% \caption{Mapping from original air-quality labels to the classes used in our experiments.}
% \label{tab:air-quality-mapping}
% \begin{tabular}{|c|l|l|r|}
% \hline
% \textbf{Class} & \textbf{Original label(s)} & \textbf{Unified name} & \textbf{Count}\\
% \hline
% 1 & Hazardous & Hazardous & 500\\
% \hline
% 2 & Poor & Poor & 1000\\
% \hline
% 3 & Moderate, Good & Qualified & 3500\\
% \hline
% \end{tabular}
% \end{table}

% Given this real data, we run the following code to get the HNP classifier based on the nultinomial logistic regression, multi-class support vector machine, and random forest:
% \begin{verbatim}
% hnp_umbrella(
%       S = train_df,
%       levels = c(0.1, 0.1),
%       tolerances = c(0.1, 0.1),
%       methods = "logistic", # could be replaced by "svm" and "randomforest"
%       hnp_split = list(c(train = 0.3, threshold = 0.7), c(train = 0.4, threshold = 0.55, error = 0.05), c(train = 0.95, error = 0.05)),
%       class_col = "Air_Quality_Label",
%       verbose = FALSE,
%       pretrained_model = predictor_fn,
%       pretrained_feature_names = feature_names,
%       pretrained_required_levels = c("1", "2", "3")
%     )
% \end{verbatim}


% We apply the package function \texttt{hnp\_boxplot\_from\_conf\_rdata()} for multinomial logistic regression (nnet::multinom) to repeat the 3-classes classification task with and without HNP for 100 times with desired control level 0.1 and violation tolerance 0.1 for under-classification errors of class $1$ and class $2$. We then visualize the under-classification error before and after applying the HNP-Umbrella algorithm.

% We report the 90th quantile of the under-classification error across repetitions. The results below show that our H--NP classifier trained based on multinomial logistic regression, multi-class support vector machine and random forest controls the under-classification error below the desired control level in more than 90\% of the repetitions, with only a slight decline in overall accuracy.

% \begin{figure}[htbp]
%   \centering
%   \includegraphics[width=0.8\linewidth]{real data application figures/air_quality_logistic.pdf}
%   \caption{comparison of the multinomial logistic classifier and its HNP classifier based on it in terms of overall accuracy and class-wise under-classification errors in air quality classification.}
%   \label{fig:air-quality-logistic}
% \end{figure}

% \begin{figure}[htbp]
%   \centering
%   \includegraphics[width=0.8\linewidth]{real data application figures/air_quality_svm.pdf}
%   \caption{comparison of the multi-class support vector machine classifier and its HNP classifier based on it in terms of overall accuracy and class-wise under-classification errors in air quality classification.}
%   \label{fig:air-quality-svm}
% \end{figure}

% \begin{figure}[htbp]
%   \centering
%   \includegraphics[width=0.8\linewidth]{real data application figures/air_quality_rf.pdf}
%   \caption{comparison of the random forest classifier and its HNP classifier based on it in terms of overall accuracy and class-wise under-classification errors in air quality classification.}
%   \label{fig:air-quality-rf}
% \end{figure}




% We also show \texttt{hnp\_umbrella()} is able to control under-classification errors of neural network as a classifier. A multilayer perceptron is trained as the classifier function. As shown in the figure, the model was configured with an input layer whose dimension was determined by the number of predictor variables, 4 hidden layer with dimension of 128, and an output layer with dimension corresponding to the target classes. Architecturally the network comprised an input fully connected layer followed by batch normalization, ReLU activation, and dropout, then two hidden blocks with the same sequence of operations, and finally an output layer. Batch normalization was enabled throughout the hidden representation learning process, and the dropout rate was set to 0.30 to reduce overfitting. The model was trained using the Adam optimizer with a learning rate of 0.001 and a weight decay coefficient of 0.0001. Training was conducted for 200 epochs with a mini-batch size of 64, and the training samples were randomly shuffled at each epoch. The cross-entropy loss function was used to optimize the model parameters, and gradient clipping with a maximum norm of 1.0 was applied during backpropagation to improve numerical stability. Although class weights were calculated from the training data using a square-root inverse-frequency scheme, the final advanced MLP implementation employed the unweighted cross-entropy loss. Computation was executed on a CUDA-enabled GPU.

% \begin{figure}[htbp]
%     \centering
%     \begin{tikzpicture}[
%         shorten >=1pt,->,
%         draw=black!50,
%         node distance=2.3cm,
%         neuron/.style={circle, fill=black!25, minimum size=16pt, inner sep=0pt},
%         input neuron/.style={neuron, fill=green!20, draw=green!50!black},
%         hidden neuron/.style={neuron, fill=blue!20, draw=blue!50!black},
%         output neuron/.style={neuron, fill=red!20, draw=red!50!black},
%         annot/.style={text width=6em, align=center}
%     ]

%         % 1. Input Layer
%         \foreach \name / \y in {1,...,4}
%             \node[input neuron] (I-\name) at (0,-\y) {};
%         \node[annot, above of=I-1, node distance=1cm] {Input Layer \\ \footnotesize Standardized Features \\ \footnotesize (dim = 9)};
%         \node[below=0.2cm of I-4] {$\vdots$};

%         % Hidden Layer 1 (128)
%         \foreach \name / \y in {1,...,4}
%             \node[hidden neuron] (H1-\name) at (2.5,-\y+0.5) {};
%         \node[annot, above of=H1-1, node distance=1cm] {Dense (128) \\ \footnotesize ReLU};
%         \node[below=0.2cm of H1-4] {$\vdots$};

%         % Hidden Layer 2 (128)
%         \foreach \name / \y in {1,...,4}
%             \node[hidden neuron] (H2-\name) at (5,-\y+0.5) {};
%         \node[annot, above of=H2-1, node distance=1cm] {Dense (128) \\ \footnotesize ReLU};
%         \node[below=0.2cm of H2-4] {$\vdots$};

%         % Hidden Layer 3 (128)
%         \foreach \name / \y in {1,...,4}
%             \node[hidden neuron] (H3-\name) at (7.5,-\y+0.5) {};
%         \node[annot, above of=H3-1, node distance=1cm] {Dense (128) \\ \footnotesize ReLU};
%         \node[below=0.2cm of H3-4] {$\vdots$};

%         % Hidden Layer 4 (128)
%         \foreach \name / \y in {1,...,4}
%             \node[hidden neuron] (H4-\name) at (10,-\y+0.5) {};
%         \node[annot, above of=H4-1, node distance=1cm] {Dense (128) \\ \footnotesize ReLU};
%         \node[below=0.2cm of H4-4] {$\vdots$};

%         % Output Layer (3 Classes)
%         \foreach \name / \y in {1,...,3}
%             \node[output neuron] (O-\name) at (12.5,-\y-0.2) {};
%         \node[annot, above of=O-1, node distance=1.3cm] {Output Layer \\ \footnotesize 3 Classes \\ \footnotesize Softmax};

%         % Connections
%         \foreach \source in {1,...,4}
%             \foreach \dest in {1,...,4}
%                 \draw (I-\source) -- (H1-\dest);

%         \foreach \source in {1,...,4}
%             \foreach \dest in {1,...,4}
%                 \draw (H1-\source) -- (H2-\dest);

%         \foreach \source in {1,...,4}
%             \foreach \dest in {1,...,4}
%                 \draw (H2-\source) -- (H3-\dest);

%         \foreach \source in {1,...,4}
%             \foreach \dest in {1,...,4}
%                 \draw (H3-\source) -- (H4-\dest);

%         \foreach \source in {1,...,4}
%             \foreach \dest in {1,...,3}
%                 \draw (H4-\source) -- (O-\dest);

%     \end{tikzpicture}
%     \caption{MLP Architecture for Air Quality Multi-class Classification (4 Hidden Layers, 128 Units Each)}
%     \label{fig:mlp_arch}
% \end{figure}

% Given the real data of air quality, we run the following code to get the HNP classifier bsed on the multilayer perceptron:
% \begin{verbatim}
% hnp_umbrella( # nolint
%       S = train_df,
%       levels = c(0.1, 0.1),
%       tolerances = c(0.1, 0.15),
%       hnp_split = list(c(train = 0.3, threshold = 0.7), c(train = 0.4, threshold = 0.55, error = 0.05), c(train = 0.95, error = 0.05)),
%       class_col = "Air_Quality_Label",
%       verbose = FALSE,
%       pretrained_model = predictor_fn,
%       pretrained_feature_names = feature_names,
%       pretrained_required_levels = c("1", "2", "3")
%     )
% \end{verbatim}


% We apply the package function \texttt{hnp\_box\_plot\_nn\_cached()} for mltilayer perceptron to repeat the 3-classes classification task with and without HNP for 100 times with desired control level 0.05 and violation tolerance 0.1 for under-classification errors of class $1$, and desired control level 0.1 and violation tolerance 0.2 for under-classification errors of class $2$. We then visualize the under-classification error before and after applying the HNP-Umbrella algorithm.

% We report the 90th quantile of the under-classification error across repetitions. Similarly, the results below show that our H--NP classifier based on multilayer perceptron controls the under-classification error of class $1$ below the desired control level 0.05 in more than 95\% of the repetitions and controls the under-classification error of class $2$ below the desired control level 0.1 in more than 80\% of the repetitions, with the cost of increasing underclassification error of class 2 and slight increasing in overall error.

% \begin{figure}[htbp]
%   \centering
%   \includegraphics[width=0.8\linewidth]{real data application figures/nn.pdf}
%   \caption{comparison of the multilayer perceptron and its HNP classifier based on it in terms of overall accuracy and class-wise under-classification errors in air quality classification.}
%   \label{fig:air-quality-nn}
% \end{figure}





In this section, we present another simulation setting for 5-class classification, similar to the one used in Simulation 2. The only modification is that the control level $\alpha_i$'s and tolerance $\delta_i$'s are all set to $0.05$. The results are summarized in Table \ref{tab:five_class_base_methods_005} and Figure \ref{fig:five_class_base_methods_boxplots}. This simulation again confirms the high-probability control of under-classification errors by H--NP classifiers, while the increase in overall classification error remains mild.













\begin{table}[h!]
\centering
\resizebox{\textwidth}{!}{%
\begin{tabular}{c !{\vrule width 1.2pt} c !{\vrule width 1.2pt}  c c c c !{\vrule width 1.2pt} c c c c !{\vrule width 1.2pt} c}
\Xhline{1.2pt}
Base method & Paradigm 
& $R_{1\star}$ & $R_{2\star}$ & $R_{3\star}$ & $R_{4\star}$ 
& $V_1$ & $V_2$ 
& $V_3$ & $V_4$ & $R_{\text{overall}}$\\
\Xhline{1.2pt}

\multirow{2}{*}{Logistic regression}
& Classical 
& 0.491 & 0.641 & 0.113 & 0.152
& 1.000 & 1.000 & 1.000 & 1.000 & 0.538 \\
& \textcolor{black}{H--NP}  
& 0.023 & 0.023 & 0.002 & 0.006
& 0.028 & 0.017 & 0.000 & 0.000 & 0.598 \\

\Xhline{1.2pt}

\multirow{2}{*}{Random forest}
& Classical
& 0.425 & 0.313 & 0.196 & 0.147
& 1.000 & 1.000 & 1.000 & 1.000 & 0.474 \\
& \textcolor{black}{H--NP}  
& 0.025 & 0.026 & 0.016 & 0.011
& 0.015 & 0.025 & 0.001 & 0.000 & 0.567 \\

\Xhline{1.2pt}

\multirow{2}{*}{SVM}
& Classical
& 0.307 & 0.143 & 0.142 & 0.145
& 1.000 & 1.000 & 1.000 & 1.000 & 0.428 \\
& \textcolor{black}{H--NP}  
& 0.024 & 0.027 & 0.021 & 0.016
& 0.031 & 0.032 & 0.002 & 0.000 & 0.519 \\

\Xhline{1.2pt}
\end{tabular}
}
\caption{Average overall classification error $R_{\text{overall}}$, average under-classification errors $R_{i\star}$, and violation rates $V_i$ for $i \in [4]$ in Section \ref{sec:extra-sim}, with $\alpha_i = \delta_i = 0.05$ for all $i \in [4]$.}
\label{tab:five_class_base_methods_005}
\end{table}



\begin{figure}[h!]
\centering

\begin{subfigure}{0.95\textwidth}
    \centering
    \includegraphics[width=\textwidth]{\detokenize{randomforest005five_classes.pdf}}
    \caption{Random forest}
\end{subfigure}

\vspace{0.5em}

\begin{subfigure}{0.95\textwidth}
    \centering
    \includegraphics[width=\textwidth]{\detokenize{logistic005five_classes.pdf}}
    \caption{Logistic regression}
\end{subfigure}

\vspace{0.5em}

\begin{subfigure}{0.95\textwidth}
    \centering
    \includegraphics[width=\textwidth]{\detokenize{svm005five_classes.pdf}}
    \caption{SVM}
\end{subfigure}

\caption{Boxplots of under-classification errors and overall classification errors for different base methods in Section \ref{sec:extra-sim}, with $\alpha_i=\delta_i=0.05$ for all $i\in[4]$.}
\label{fig:five_class_base_methods_boxplots}
\end{figure}





\FloatBarrier
